\documentclass[11pt]{article}

% ---------------------------------------------------------------------------
% Docs/theory.tex -- self-contained theoretical companion for the HUM (Hilbert
% Uniqueness Method) approach to controllability of the 1D heat equation.
%
% Source: Proposal/THEORY_EXTRACTION.md (the signed-off extraction of theory
% from Cipriano Callejas Hernandez's 2024 CIMAT MSc thesis, "Numerical and
% constrained controllability of the heat equation", plus independently-cited
% standard results explicitly marked as such throughout). Internal
% (distributed) control is treated as the primary case; boundary control
% appears only as an explicit point of contrast, per that document's scope.
%
% Build: `make` in this directory (see Makefile), or directly:
%   latexmk -pdf theory.tex
% ---------------------------------------------------------------------------

\usepackage[T1]{fontenc}
\usepackage[utf8]{inputenc}
\usepackage[margin=1in]{geometry}
\usepackage{amsmath, amssymb, amsthm, mathtools}
\usepackage{enumitem}
\usepackage{xcolor}
\usepackage[colorlinks=true, linkcolor=blue!50!black, citecolor=blue!50!black,
            urlcolor=blue!50!black]{hyperref}
\usepackage[round]{natbib}

% --- Theorem-like environments -------------------------------------------
\theoremstyle{plain}
\newtheorem{theorem}{Theorem}[section]
\newtheorem{proposition}[theorem]{Proposition}
\newtheorem{lemma}[theorem]{Lemma}
\newtheorem{corollary}[theorem]{Corollary}
\theoremstyle{definition}
\newtheorem{definition}[theorem]{Definition}
\theoremstyle{remark}
\newtheorem{remark}[theorem]{Remark}

% --- Convenience macros -----------------------------------------------------
\newcommand{\R}{\mathbb{R}}
\DeclareMathOperator{\rank}{rank}

\title{The Hilbert Uniqueness Method for Internal Control of the\\
  One-Dimensional Heat Equation: A Theoretical Companion}
\author{Compiled from Cipriano Callejas Hern\'andez's CIMAT MSc thesis\\
  and independently-cited standard results (see text)}
\date{}

\begin{document}
\maketitle

\begin{abstract}
\noindent
This document collects, with precise statements and citations, the continuous
and discrete controllability theory underlying the Hilbert Uniqueness Method
(HUM) applied to \emph{internal} (distributed) control of the one-dimensional
heat equation, together with the finite-dimensional (ODE) controllability
results that justify passing from the continuous problem to a semi-discrete
one. Boundary control is discussed only where it clarifies a genuine point of
contrast with the internal case. Results are drawn from Callejas
Hern\'andez's 2024 CIMAT thesis, \emph{Numerical and constrained
controllability of the heat equation}, whenever available; every result
\emph{not} found there is cited independently and marked as such, never
attributed to the thesis. Open questions and known limitations are stated
explicitly rather than smoothed over.
\end{abstract}

\tableofcontents

\section{Well-posedness of the heat equation}
\label{sec:wellposedness}

\subsection*{Abstract standing assumption}

Let $T>0$ be the (fixed) control time horizon, and let $H$ and $U$ be real
Hilbert spaces -- the \emph{state space} and \emph{control space},
respectively (specialized to concrete function spaces below, e.g.\
$H=U=L^2(0,L)$ for internal control, where $L>0$ is the length of the spatial
interval). Controllability problems are then posed abstractly as
\begin{equation}
  \label{eq:abstract-system}
  y'(t) + Ay(t) = Bu(t) \ \text{in } (0,T), \qquad y(0) = y_0 \in H,
\end{equation}
where $u \in L^2(0,T;U)$ is the control, and $A, B$ are (possibly unbounded)
linear operators densely defined on $D(A) \subset H$ and $D(B) \subset U$. It
is assumed throughout that \eqref{eq:abstract-system} is \emph{well posed}: for
every $(y_0, u) \in H \times L^2(0,T;U)$ there is a unique solution
$y \in C([0,T];H)$ depending continuously on the data, understood in the weak
sense via the semigroup generated by $-A$ \citep{pazy2012semigroups}.

\subsection{Internal (distributed) control}
\label{ssec:wp-internal}

For internal control, $H = U = L^2(0,L)$, $A = -\alpha\partial_{xx}$ for a
fixed diffusion coefficient $\alpha>0$, with
$D(A) = H^2(0,L)\cap H^1_0(0,L)$, and $B$ is the (bounded) multiplication
operator by $\mathbf{1}_\omega$ for a subset $\omega \Subset (0,L)$:
\begin{equation}
  \label{eq:heat-internal}
  \begin{cases}
    y_t - \alpha y_{xx} = v(x,t)\mathbf{1}_\omega(x), & (x,t) \in (0,L)\times(0,T),\\
    y(0,t) = y(L,t) = 0, & t \in (0,T),\\
    y(x,0) = y_0(x), & x \in (0,L).
  \end{cases}
\end{equation}

\begin{theorem}[Well-posedness, internal control]
  \label{thm:wp-internal}
  For every $y_0 \in L^2(0,L)$ and $v \in L^2\big((0,L)\times(0,T)\big)$,
  problem \eqref{eq:heat-internal} has a unique solution
  $y \in C\big([0,T];L^2(0,L)\big)$. Moreover, \eqref{eq:heat-internal} is null
  controllable \citep{fursikov1996imanuvilov}.
\end{theorem}

\subsection{Boundary control}
\label{ssec:wp-boundary}

For boundary control, the same equation is instead driven through a Dirichlet
datum at one endpoint:
\begin{equation}
  \label{eq:heat-boundary}
  \begin{cases}
    y_t - \alpha y_{xx} = 0, & (x,t) \in (0,L)\times(0,T),\\
    y(0,t) = 0, \quad y(L,t) = v(t), & t \in (0,T),\\
    y(x,0) = y_0(x), & x \in (0,L).
  \end{cases}
\end{equation}
Because the control enters through the boundary trace rather than as a bounded
source term, \eqref{eq:heat-boundary} only admits a solution \emph{by
transposition}, and in a weaker space than the internal case:

\begin{theorem}[Well-posedness, boundary control]
  \label{thm:wp-boundary}
  For every $y_0 \in L^2(0,L)$ and $v \in L^2(0,T)$, problem
  \eqref{eq:heat-boundary} has a unique weak (transposition) solution
  $y \in C\big([0,T]; H^{-1}(0,L)\big)$, and \eqref{eq:heat-boundary} is null
  controllable \citep{fursikov1996imanuvilov}.
\end{theorem}

From here on, $\Omega$ denotes the spatial domain of a result stated at the
general (possibly multi-dimensional) level of the underlying PDE-control
literature; in the concrete one-dimensional setting of
\eqref{eq:heat-internal}--\eqref{eq:heat-boundary} used throughout the rest of
this document, $\Omega = (0,L)$. Correspondingly, $Q := \Omega\times(0,T)$ is
the space-time cylinder and $\Sigma := \partial\Omega\times(0,T)$ is its
lateral boundary (where the boundary control $v(t)$ of
\eqref{eq:heat-boundary} acts).

\begin{remark}
  The objectives of exact/null controllability require $y(\cdot,T) \in
  L^2(\Omega)$, a strictly stronger requirement than the ambient
  well-posedness class $y \in C([0,T];H^{-1}(\Omega))$ \citep{belgacem2011dirichlet}.
  This gap between the natural well-posedness space and the space in which
  controllability is sought is a recurring theme (see the ill-posedness
  discussion in Section~\ref{sec:known-issues}).
\end{remark}

\subsubsection*{Solution-by-transposition construction}

The rigorous justification for Theorem~\ref{thm:wp-boundary} splits the
solution as $y = z + w$, where $w$ solves an auxiliary problem with the same
data but homogeneous boundary conditions, and $z$ is defined by transposition
against a smooth adjoint problem.

\begin{proposition}
  \label{prop:transposition}
  The auxiliary solution $w$ satisfies
  $w \in L^2(0,T;L^2(\Omega)) \cap C^0([0,T];H^{-1}(\Omega))$
  \citep{evans2010partial}. The adjoint problem
  \[
    -p_t - \Delta p = f \ \text{in } Q, \qquad p = 0 \ \text{on } \Sigma,
    \qquad p(x,T) = 0 \ \text{in } \Omega,
  \]
  with $f \in L^2(0,T;L^2(\Omega))$, has a unique solution
  $p \in L^2\big(0,T; H^1_0(\Omega)\cap H^2(\Omega)\big)$ with
  $p_t \in L^2(Q)$; the transposition identity built from $p$ then yields
  $z \in C^0([0,T];H^{-1}(\Omega))$, so $y = z+w$ is well defined
  \citep{lions2012non}.
\end{proposition}

This construction is also the continuous-theory root of any numerical HUM
variant for boundary control that works directly with the (rougher)
transposition solution rather than with the adjoint state.

\begin{remark}[Gap in the thesis, now filled independently -- not from the
  thesis]
  \label{rem:energy-wellposedness}
  No statement of well-posedness in the classical parabolic energy form
  $y \in C([0,T];L^2(\Omega)) \cap L^2(0,T;H^1_0(\Omega))$ appears in the
  thesis; it consistently works either in the abstract semigroup framework
  above, or in the transposition/weak-solution framework in $H^{-1}$. For
  completeness, the standard energy estimate is:
  \begin{theorem}[Classical energy estimate]
    For $y_0 \in L^2(\Omega)$ and $f \in L^2(0,T;L^2(\Omega))$, the linear heat
    equation $y_t - \Delta y = f$ in $Q$, $y=0$ on $\Sigma$, $y(\cdot,0)=y_0$,
    has a unique weak solution
    $y \in C([0,T];L^2(\Omega)) \cap L^2(0,T;H^1_0(\Omega))$, with
    $y_t \in L^2(0,T;H^{-1}(\Omega))$, obtained via the Galerkin method and the
    energy estimate
    \[
      \sup_{t\in[0,T]}\|y(t)\|_{L^2(\Omega)}^2
      + \int_0^T \|y(t)\|_{H^1_0(\Omega)}^2\,dt
      \;\le\; C\Big(\|y_0\|_{L^2(\Omega)}^2 + \|f\|_{L^2(Q)}^2\Big).
    \]
  \end{theorem}
  \citep[\S 7.1]{evans2010partial}. This is the same book already cited in
  Proposition~\ref{prop:transposition} for a different, thesis-attributed
  result; unlike every thesis-sourced statement in this document (checked
  line-by-line against the primary source), this statement was not verified
  page-by-page against Evans' book -- it is standard textbook material, not a
  literal quote.
\end{remark}

\section{Regularity of states, controls, and the adjoint state}
\label{sec:regularity}

The thesis states no separate, numbered "regularity" theorem for the
HUM-optimal control or the state; what follows is what genuinely exists,
plus (marked explicitly) what is added independently to complete the picture.

\subsection{Internal control}

The adjoint final datum lives in $\varphi_T \in L^2(0,1)$ -- no further
regularity statement beyond the well-posedness class of
Section~\ref{ssec:wp-internal} is given.

\subsection{Boundary control}

The boundary case's adjoint final datum lives in a strictly smaller space,
flagged in the thesis as the key point of departure from the internal case:

\begin{remark}
  \label{rem:boundary-h10}
  For boundary control, the adjoint system has final datum
  $\varphi_T \in H^1_0(0,1)$, rather than $L^2(0,1)$ as in the internal case.
  This single distinction is the source of the $H^1_0$/$H^{-1}$ machinery
  needed throughout the boundary theory (Section~\ref{ssec:wp-boundary},
  Proposition~\ref{prop:transposition}).
\end{remark}

\subsection{Shared / adjacent}

The construction of higher-regularity controls (Appendix~A.2 of the thesis) is
technical machinery specific to the \emph{constrained} (state-positivity)
problem of Section~\ref{sec:constrained-aside}, not a general HUM
control-regularity result for the unconstrained baseline treated here -- see
that section for the statement (Proposition~\ref{prop:hidden-regularity}).
Separately, control regularity for constrained minimal-time controllability is
addressed in \citet[Thm.~4]{loheac2017minimal}, again in the constrained
setting rather than the baseline.

\begin{remark}[Gap in the thesis, now filled independently -- not from the
  thesis]
  \label{rem:smoothing}
  No parabolic-smoothing theorem is stated in the thesis, though it is
  implicit in $D(A) = H^2(\Omega)\cap H^1_0(\Omega)$. Since $A=-\Delta$
  generates an analytic semigroup $(e^{-tA})_{t\ge0}$ on $L^2(\Omega)$
  \citep{pazy2012semigroups}, the uncontrolled solution satisfies the standard
  regularizing effect
  \[
    y_0 \in L^2(\Omega) \ \Longrightarrow\
    y(t) \in D(A) \subset H^1_0(\Omega)\cap H^2(\Omega) \ \text{for every } t>0,
    \qquad \|y(t)\|_{D(A)} \le \frac{C}{t}\|y_0\|_{L^2(\Omega)}.
  \]
  As with Remark~\ref{rem:energy-wellposedness}, this citation is the same
  book already used for the abstract semigroup framework, and this specific
  statement was not checked page-by-page against it.
\end{remark}

\begin{remark}
  The finite-dimensional Kalman-condition/observability machinery underlying
  \emph{discrete} uniform controllability (Section~\ref{sec:discrete}) is a
  separate notion -- ``regularity of controllability'' for a finite-dimensional
  ODE system -- distinct from the PDE state/control regularity discussed here,
  and is deliberately not conflated with it.
\end{remark}

\section{Controllability definitions}
\label{sec:definitions}

These definitions are shared across internal and boundary control -- the
thesis does not split them, and neither do we. Here $H$ and $U$ are the state
and control spaces of \eqref{eq:abstract-system}, as fixed in
Section~\ref{sec:wellposedness}.

\begin{definition}[Controllability notions for \eqref{eq:abstract-system}]
  \label{def:controllability}
  Let $T > 0$. System \eqref{eq:abstract-system} is said to be:
  \begin{enumerate}[label=(\roman*)]
    \item \textbf{exactly controllable} in time $T$ if, for all
      $(y_0,y_1)\in H\times H$, there exists $u \in L^2(0,T;U)$ such that the
      solution satisfies $y(T;y_0,u) = y_1$;
    \item \textbf{exactly controllable to trajectories} in time $T$ if, for
      all $(y_0,\hat y_0)\in H\times H$ and $\hat u \in L^2(0,T;U)$, there
      exists $u \in L^2(0,T;U)$ such that
      $y(T;y_0,u) = y(T;\hat y_0,\hat u)$;
    \item \textbf{null controllable} (or exactly controllable to zero) in
      time $T$ if, for all $y_0 \in H$, there exists $u \in L^2(0,T;U)$ such
      that $y(T;y_0,u)=0$;
    \item \textbf{approximately controllable} in time $T$ if, for all
      $(y_0,y_1)\in H\times H$ and every $\varepsilon>0$, there exists
      $u \in L^2(0,T;U)$ such that $\|y(T;y_0,u)-y_1\|_H < \varepsilon$.
  \end{enumerate}
\end{definition}

Equivalently, defining the reachable-set map
$\Phi: L^2(0,T;U)\to H$, $u \mapsto y(T;y_0,u)$, and its range
$R(T) = \{y(T;y_0,u): u\in L^2(0,T;U)\} = \Phi\big(L^2(0,T;U)\big) \subset H$,
the four notions above correspond to
$R(T)=H$, $R(T) \ni y(T;\hat y_0,\hat u)$ for every choice, $0 \in R(T)$, and
$R(T)$ dense in $H$, respectively. ($\Phi$ is distinct from the completion
space $F$ introduced in Remark~\ref{rem:hum-existence} below -- different
symbol, different role.)

\begin{remark}
  When there is no restriction on the norm of the initial condition, the
  properties in Definition~\ref{def:controllability} hold \emph{globally};
  under such a restriction, controllability holds only \emph{locally}.
\end{remark}

\subsection*{Steady-state controllability (constrained-problem aside)}

For the constrained/steady-state problem of Section~\ref{sec:constrained-aside}
the thesis restates controllability more concretely for the
boundary-controlled equation and adds:

\begin{definition}
  System \eqref{eq:heat-boundary} is (exactly) controllable to
  $y_1 \in L^2(\Omega)$ if, for every $y_0 \in L^2(\Omega)$, there exist
  $T_{y_0}>0$ and $u \in L^2(\Sigma)$ such that the solution satisfies
  $y(T_{y_0}) = y_1$.
\end{definition}

\begin{definition}[Steady state]
  $\bar y$ is a steady state of \eqref{eq:heat-boundary} if $-\Delta\bar y = 0$
  in $\Omega$, $\bar y = \bar u$ on $\partial\Omega$.
\end{definition}

\begin{definition}[Steady-state controllability]
  Let $y_0,y_1$ be steady states. System \eqref{eq:heat-boundary} is
  exactly-steady-state controllable in time $T>0$ if there exists
  $u \in L^2(\Sigma)$ such that the solution satisfies $y(\cdot,0)=y_0$ and
  $y(\cdot,T)=y_1$.
\end{definition}

\section{HUM duality theory}
\label{sec:hum-duality}

\subsection{The finite-dimensional (ODE) prototype}
\label{ssec:ode-prototype}

The infinite-dimensional theory below generalizes a finite-dimensional
template for the linear ODE system
\begin{equation}
  \label{eq:ode-system}
  y'(t) = Ay(t) + Bv(t) \ \text{in } (0,T), \qquad y(0) = y_0,
\end{equation}
with $n,m \in \mathbb{N}$ (the state and control dimensions),
$A \in \R^{n\times n}$, $B \in \R^{n\times m}$, $y : [0,T]\to\R^n$,
$v \in L^2(0,T;\R^m)$. This template is stated here in full because it is the
direct prerequisite for the discrete/ODE uniform-controllability theory of
Section~\ref{sec:discrete}: a space semi-discretization reduces the PDE
control problem to exactly this ODE setting.

\begin{theorem}[Kalman rank condition]
  \label{thm:kalman}
  Let $K = [B, AB, \dots, A^{n-1}B]$. System \eqref{eq:ode-system} is exactly
  controllable in every time $T>0$ if and only if $\rank K = n$.
\end{theorem}
\begin{proof}[Source]
  \citet[Thm.~1.16]{coron2007control}.
\end{proof}

The adjoint system of \eqref{eq:ode-system} is
\begin{equation}
  \label{eq:ode-adjoint}
  \varphi'(t) = -A^*\varphi(t) \ \text{in } (0,T), \qquad \varphi(T) = \varphi_T.
\end{equation}

\begin{theorem}[Controllability $\Leftrightarrow$ observability]
  \label{thm:ode-observability}
  System \eqref{eq:ode-system} is controllable in time $T$ if and only if the
  adjoint system \eqref{eq:ode-adjoint} is \emph{observable} in time $T$, i.e.\
  there exists $C=C(T)>0$ such that, for every solution $\varphi$ of
  \eqref{eq:ode-adjoint},
  \begin{equation}
    \label{eq:ode-observability-ineq}
    |\varphi(0)|^2_{\R^n} \le C\int_0^T |B^*\varphi(t)|^2_{\R^m}\,dt.
  \end{equation}
  Both properties hold in every time $T$ if and only if the Kalman rank
  condition (Theorem~\ref{thm:kalman}) holds.
\end{theorem}
\begin{proof}[Source]
  \citet[Thm.~2.1.1]{zuazua2002controllability}.
\end{proof}

\begin{theorem}[HUM control, finite-dimensional]
  \label{thm:ode-hum}
  Let $\mathcal{J}: \R^n \to \R$,
  $\mathcal{J}(\varphi_T) = \tfrac12\int_0^T |B^*\varphi|^2 + (y_0,\varphi(0))$,
  where $\varphi$ solves \eqref{eq:ode-adjoint} with $\varphi(T)=\varphi_T$.
  If $\mathcal{J}$ has a minimizer $\hat\varphi_T$, then, letting $\hat\varphi$
  solve \eqref{eq:ode-adjoint} with $\hat\varphi(T)=\hat\varphi_T$, the control
  $v = B^*\hat\varphi$ is the minimal-$L^2$-norm control steering
  \eqref{eq:ode-system} from $y_0$ to $0$ at time $T$.
\end{theorem}

The construction hinges on the \emph{Gramian operator}
\begin{equation}
  \label{eq:ode-gramian}
  \Lambda: \R^n \to \R^n, \qquad
  \Lambda(g) := \int_0^T B^*\varphi_g(s)\,ds,
\end{equation}
where $\varphi_g$ solves \eqref{eq:ode-adjoint} with final datum $g$; $\Lambda$
is well defined by duality \citep{boyer2013penalised}. This is the same
Gramian apparatus used below for the PDE-level construction, over $\R^n$
instead of an infinite-dimensional Hilbert space.

\subsection{Infinite-dimensional generalization}
\label{ssec:hum-infinite}

In the PDE setting the Kalman rank condition has no direct analogue; more
delicate tools are needed, and in general exact controllability of the heat
equation fails while approximate controllability holds
\citep{zuazua2002controllability} -- see also \citet{coron2007control,
zabczyk2020mathematical} for the general control-theoretic background this
document draws on. Approximate controllability of the 1D heat equation can
also be established directly via a variational approach
\citep{fabre1995approximate}, as an alternative to the HUM route developed
below. Nonetheless the HUM construction extends: with $H$, $U$ the state and
control spaces of Section~\ref{sec:wellposedness}, and adjoint system in $H$,
\begin{equation}
  \label{eq:pde-adjoint}
  \varphi'(t) = -A^*\varphi(t) \ \text{in } (0,T), \qquad \varphi(T) = g \in H,
\end{equation}

\subsubsection*{Internal control}

\begin{proposition}
  \label{prop:pde-gramian}
  The Gramian $\Lambda: H \to L^2(0,T;U)$ is a positive-definite, symmetric,
  bounded linear operator \citep{boyer2013penalised}, giving rise to the
  bilinear form and dual functional
  \begin{equation}
    \label{eq:pde-functional}
    a(g,\xi) = \int_0^T B^*\varphi_g \cdot B^*\varphi_\xi\,dt, \qquad
    J(g) = \tfrac12 a(g,g) + \langle y_0, \varphi_g(0)\rangle,
  \end{equation}
  with $\langle\cdot,\cdot\rangle$ the $H\times H'$ duality pairing. For
  internal control, $H = U = L^2(0,1)$ and $B^* = B = \mathbf{1}_\omega$.
\end{proposition}

The thesis develops \eqref{eq:pde-functional} directly and moves to
\emph{penalized} HUM (Section~\ref{sec:penalized}) as its route to a
well-posed minimization, rather than stating a standalone
existence/uniqueness theorem for the unpenalized functional.

\begin{remark}[Gap in the thesis, now filled independently -- not from the
  thesis]
  \label{rem:hum-existence}
  The classical existence/uniqueness argument for the unpenalized
  infinite-dimensional minimizer is the original Lions completion-space
  construction. Let
  $F := $ the completion of $H$ with respect to the seminorm
  $\|g\|_F := a(g,g)^{1/2}$. If this seminorm is genuinely a norm on $H$ --
  equivalent to a unique-continuation property for \eqref{eq:pde-adjoint}
  ($B^*\varphi_g \equiv 0$ on $(0,T) \Rightarrow g=0$) -- then $H$ embeds
  continuously and densely into $F$, the linear form
  $g \mapsto -\langle y_0,\varphi_g(0)\rangle$ extends continuously to $F$,
  and Riesz representation on the Hilbert space $F$ yields a unique
  $\hat g \in F$ minimizing $J$ over $F$. This formalizes what the thesis's
  Remark~\ref{rem:coercivity-space} (Section~\ref{sec:known-issues}) describes
  informally as ``the space where $J$ is coercive.''
  Citation: \citet{lions1988controlabilite} -- the original HUM source. This
  is the first citation of Lions' book \emph{directly} in this document; the
  thesis itself cites it only indirectly, via \citealp{russell1990j}'s review
  of it. As with Remarks~\ref{rem:energy-wellposedness} and~\ref{rem:smoothing},
  no specific page/theorem number is claimed here.
\end{remark}

\begin{theorem}[Observability inequality, internal control]
  \label{thm:internal-observability}
  If there exists $C>0$ such that
  \begin{equation}
    \label{eq:internal-observability}
    |\varphi(\cdot,0)|^2_{L^2(0,1)} \le
    C \int_0^1\!\int_0^T |\mathbf{1}_\omega\varphi|^2\,dx\,dt,
  \end{equation}
  then $J$ is coercive.
\end{theorem}

\begin{remark}
  The proof of \eqref{eq:internal-observability} itself requires Carleman
  estimates and is deferred to \citet{fursikov1996imanuvilov}.
\end{remark}

\begin{remark}[Gap in the thesis, now filled independently -- not from the
  thesis: an alternative route]
  \label{rem:lebeau-robbiano}
  Carleman estimates are not the only route to
  \eqref{eq:internal-observability}. The Lebeau--Robbiano spectral inequality
  \[
    \|u\|^2_{L^2(\Omega)} \le C e^{C\lambda}\|u\|^2_{L^2(\omega)}
  \]
  -- valid for any finite combination $u$ of Dirichlet-Laplacian
  eigenfunctions with eigenvalues $\le\lambda$, and any open $\omega\subset
  \Omega$ -- yields null controllability of the heat equation from any open
  subset, in any time $T>0$, via an iterative/telescoping argument
  \citep{lebeau1995controle}. This citation, like
  Remarks~\ref{rem:energy-wellposedness}, \ref{rem:smoothing}
  and~\ref{rem:hum-existence}, is standard-form and was not verified
  page-by-page against the original paper.
\end{remark}

\subsubsection*{Boundary control}

Per Remark~\ref{rem:boundary-h10}, the boundary dual functional
$J_b(\varphi_T)$ lives over $\varphi_T \in H^1_0(0,1)$ rather than
$L^2(0,1)$. No further boundary-specific infinite-dimensional HUM theorem is
stated beyond the general development above.

\section{Penalized HUM: convergence and threshold results}
\label{sec:penalized}

To obtain a well-posed minimization even when the unpenalized functional $J$
of \eqref{eq:pde-functional} is not coercive on $H$, one instead minimizes the
penalized functional
\begin{equation}
  \label{eq:penalized-functional}
  J_\varepsilon(g) = \tfrac12 a(g,g) + \langle y_0,\varphi_g(0)\rangle
    + \varepsilon\|g\|_H, \qquad \varepsilon>0.
\end{equation}

\begin{proposition}[Critical-point inequality]
  \label{prop:critical-point}
  If $J_\varepsilon$ attains its minimum at $\hat g \in H$, then for every
  $\xi \in H$,
  \[
    a(\hat g,\xi) + \langle y_0,\varphi_\xi(0)\rangle \le \varepsilon\|\xi\|_H,
  \]
  and consequently the controlled solution satisfies
  $\|y_{\hat v}(T)\| \le \varepsilon$, where $\hat v = B^*\hat\varphi$.
\end{proposition}

\begin{proposition}[Penalized-HUM duality]
  \label{prop:penalized-duality}
  Let $F_\varepsilon(v) := F_1(v) + \tfrac{1}{2\varepsilon}\|y_v(T)\|_H$
  (the primal functional). Then
  $F_\varepsilon(\hat v) = -J_\varepsilon(\hat g)$, a consequence of the
  Fenchel--Rockafellar duality theorem \citep{ekeland1999convex} that also
  holds for $\varepsilon=0$. Moreover $y(T;y_0,\hat v) = -\varepsilon\hat g$,
  $\|y(T;y_0,\hat v)\|_H \le \|y(T;y_0,0)\|_H$, and
  $\|\hat g\|_H \le \tfrac1\varepsilon\|y(T;y_0,0)\|_H$.
\end{proposition}
\begin{proof}[Source]
  \citet[Prop.~1.5]{boyer2013penalised}.
\end{proof}

\begin{theorem}[Convergence as $\varepsilon\to0$]
  \label{thm:penalized-convergence}
  Let $y_\varepsilon(T)$ be the state reached by the penalized-HUM control
  $v_\varepsilon = \hat v(\varepsilon)$. Then
  \[
    y_\varepsilon(T) \to \mathbb{P}_{Q_F}\big(y_f(T)\big) \ \text{as } \varepsilon\to0,
  \]
  where $Q_F := \{g : B^*\varphi_g \equiv 0 \text{ on } (0,T)\} = \ker\Lambda$
  is the set of non-observable states, $\mathbb{P}_{Q_F}$ is the orthogonal
  projection onto $Q_F$, and $y_f = y(\cdot\,;y_0,0)$ is the free solution. In
  particular, if $Q_F = \{0\}$ then $y_\varepsilon(T) \to 0$.
\end{theorem}
\begin{proof}[Source]
  \citet[Thm.~1.11]{boyer2013penalised}.
\end{proof}

\begin{proposition}
  \label{prop:penalized-strong-conv}
  \begin{enumerate}[label=(\roman*)]
    \item If the primal system is approximately or null controllable, then
      $y(T;y_0,\hat v_\varepsilon) \to 0$ as $\varepsilon \to 0$.
    \item If the primal system is null controllable, then
      $v_\varepsilon \to v^0$ strongly in $L^2(0,T;U)$ as $\varepsilon\to0$,
      where $v^0$ is the minimal-norm null control (null controllability is
      needed for the strong convergence, since it guarantees
      $\inf_{L^2(0,T;U)}F_\varepsilon < \infty$ as $\varepsilon\to0$).
  \end{enumerate}
\end{proposition}
\begin{proof}[Source]
  \citet[Prop.~1.7]{boyer2013penalised}.
\end{proof}

\begin{remark}
  In finite dimension, $Q_F$ is the kernel of the adjoint Kalman matrix $K^*$
  of Theorem~\ref{thm:kalman}; hence $Q_F=\{0\}$ if and only if the Kalman
  rank condition holds.
\end{remark}

\subsection{Internal control}

Propositions~\ref{prop:pde-gramian}--\ref{prop:penalized-strong-conv} apply
directly to internal control over $H=L^2(0,1)$; nothing further is added
beyond this general machinery.

\subsection{Boundary control}

The same machinery restates over $H^1_0/H^{-1}$ per Remark~\ref{rem:boundary-h10}.
The condition-number result below is stated once, generically, and applies to
whichever Gramian is in play (internal or boundary).

\begin{proposition}[Condition number of the penalized Gramian]
  \label{prop:condition-number}
  With $\Lambda_\varepsilon := \Lambda + \varepsilon I$,
  $\|\Lambda_\varepsilon\| = \varepsilon + \|\Lambda\|$ and
  $\|\Lambda_\varepsilon^{-1}\| = \varepsilon^{-1}$, so the condition number of
  $\Lambda_\varepsilon$ is
  \[
    \nu_a = 1 + \frac{1}{\varepsilon}\|\Lambda\|.
  \]
\end{proposition}

\begin{theorem}[Conjugate-gradient convergence rate]
  \label{thm:cg-rate}
  The conjugate-gradient iterates $\{u_n\}_{n\ge1}$ for minimizing the
  bilinear form $a(\cdot,\cdot)$ satisfy
  \[
    \|u_n - u\| \le C\|u_0-u\|
      \left(\frac{\sqrt{\nu_a}-1}{\sqrt{\nu_a}+1}\right)^n, \qquad n\ge1,
  \]
  where $\nu_a = \|\Lambda\|\|\Lambda^{-1}\|$ is the condition number of
  $a(\cdot,\cdot)$: convergence is faster the closer $u_0$ is to $u$, and the
  closer $\sqrt{\nu_a}$ is to $1$.
\end{theorem}
\begin{proof}[Source]
  \citet{glowinski2008lectures}.
\end{proof}

\begin{remark}
  No further quantitative threshold or error-rate result exists beyond
  Proposition~\ref{prop:condition-number} and Theorem~\ref{thm:cg-rate}. A
  scaling law $\varepsilon \sim Ch^p$ (for a mesh-refinement parameter $h$,
  diffusion coefficient $\alpha$, and $p$ larger than the scheme's order of
  accuracy) is reported, but only as an \textbf{experimental conjecture} from
  numerical observation, not a proven theorem, and is labeled as such
  wherever referenced.
\end{remark}

\section{Discrete-in-space uniform controllability: the ODE-controllability bridge}
\label{sec:discrete}

This section covers the theory of \emph{space} semi-discretization -- the
technique actually used in this work -- and its relation to the
finite-dimensional (ODE) controllability theory of
Section~\ref{ssec:ode-prototype}. Internal (distributed) control is treated
as the primary case; boundary control appears only as a one-line contrast.

\subsection{Framing}
\label{ssec:discrete-framing}

\begin{quote}
  \itshape ``A space semi-discrete approximation of the parabolic control
  problems [\eqref{eq:heat-boundary} or \eqref{eq:heat-internal}] reduces the
  initial problem to a control problem for an ordinary differential equations
  (ODE's) system.''
\end{quote}
This reduction is the central bridge exploited throughout this section: study
of the semi-discrete heat equation's controllability becomes an application of
Theorems~\ref{thm:kalman}--\ref{thm:ode-hum} to a specific, mesh-dependent
family of linear ODE systems.

The numerical-HUM approach itself originates with the pioneering work of
\citet{glowinski1990numerical} for the wave equation, treated as an
unconstrained optimization problem for the adjoint system via convex-duality
theory \citep{ekeland1999convex}; it was later adapted to the heat equation to
find null controls by \citet{carthel1994exact}. HUM belongs to a broader
family of dual methods, including the \emph{weighted} HUM, which incorporates
a small perturbation to damp oscillatory control behavior
\citep{zuazua2006control}; more recently, primal methods working directly
with the primal system via Carleman weights \citep{fernandez2014numerical}
and least-squares/variational methods \citep{munch2014numerical} have given
results broadly comparable to the dual approach.

Discrete controllability of parabolic problems is usually studied via three
formulations -- time-discrete, space-discrete, and fully discrete
\citep{boyer2020carleman}. Uniform controllability (defined precisely in
Definition~\ref{def:uniform-controllability} below) is established
\emph{only} for the space-discrete formulation of the one-dimensional heat
equation \citep{zuazua2006control,lopez1998some}; the other two formulations
are not pursued in this work (see Section~\ref{ssec:other-formulations} for
why).

\subsection{The space-discrete ODE system}
\label{ssec:space-discrete-system}

Let $N \in \mathbb{N}$, let $h = L/(N+1)$ be the corresponding uniform mesh
size on $(0,L)$, $x_j = jh$ for $j=0,\dots,N+1$, and let
$y_h(t) = (y_1(t),\dots,y_N(t)) \in \R^N$ approximate
$\big(y(x_1,t),\dots,y(x_N,t)\big)$ via the standard central-difference
approximation of $\partial_{xx}$. For internal control, \eqref{eq:heat-internal}
becomes the linear ODE system on $\R^N$
\begin{equation}
  \label{eq:discrete-internal}
  y_h'(t) = A_h y_h(t) + B_h v_h(t), \qquad
  A_h = \frac{\alpha}{h^2}
  \begin{pmatrix}
    -2 & 1 & & \\
    1 & -2 & 1 & \\
    & \ddots & \ddots & \ddots \\
    & & 1 & -2
  \end{pmatrix} \in \R^{N\times N},
\end{equation}
where $v_h(t) \in \R^N$ and $B_h = \mathbf{1}_{\omega_h} \in \R^{N\times N}$ is
the diagonal matrix with $(\mathbf{1}_{\omega_h})_{jj}=1$ if $x_j \in \omega$
and $0$ otherwise (the identity matrix if $\omega$ is the whole domain). For
boundary control, the analogous system \citep[cf.][]{leveque2007finite} has
the same $A_h \in \R^{N\times N}$ but $v_h(t) \in \R$ a scalar and
$B_h \in \R^{N\times 1}$ a column vector, rather than $v_h(t) \in \R^N$ and
$B_h \in \R^{N\times N}$ as in \eqref{eq:discrete-internal}.

\begin{proposition}[Spectrum of $A_h$]
  \label{prop:spectrum}
  $A_h$ has eigenvalues
  $\lambda_j = \dfrac{2\alpha}{h^2}\big(1-\cos(j\pi h)\big)$, $j=1,\dots,N$,
  with eigenvectors $e^j_i = \sin(ij\pi h)$, $i=1,\dots,N$
  \citep[\S3.4]{leveque2007finite}.
\end{proposition}

This explicit, closed-form spectrum is precisely what makes the uniform
controllability results below tractable: the proof of
Theorem~\ref{thm:discrete-null-control} "relies solely on the fact that the
spectrum of the Laplacian can be computed explicitly."

\subsection{Uniform controllability}
\label{ssec:uniform-controllability}

\begin{definition}[Uniform controllability]
  \label{def:uniform-controllability}
  Controllability properties must be established via an observability
  inequality for the discrete adjoint system, rather than by directly
  invoking the Kalman rank condition, since the latter is a
  \textbf{qualitative, binary} criterion for a fixed finite-dimensional
  system: it says nothing about how the observability constant behaves as the
  system size grows with $h \to 0$. System \eqref{eq:discrete-internal} (resp.\
  its boundary analogue) is said to be \textbf{uniformly controllable with
  respect to $h$} if the discrete observability inequality below holds for a
  suitable constant $C>0$ \emph{independent of $h$}.
\end{definition}

The discrete adjoint system, with $\varphi_h(t) \in \R^N$, is
\begin{equation}
  \label{eq:discrete-adjoint}
  \varphi_h'(t) = -A_h^\top\varphi_h(t) \ \text{in } (0,T), \qquad
  \varphi_h(T) = \varphi_h^T \in \R^N.
\end{equation}

\begin{theorem}[Discrete observability inequality]
  \label{thm:discrete-observability}
  For every $T>0$ there is $C(T)>0$ such that
  \[
    h\sum_{j=1}^N |\varphi_j(0)|^2 \le C \int_0^T \left|\frac{\varphi_N(t)}{h}\right|^2 dt
  \]
  for every solution $\varphi_h$ of \eqref{eq:discrete-adjoint} and every
  $h>0$.
\end{theorem}
\begin{proof}[Source]
  \citet{lopez1998some}, based on a discrete observability inequality that
  relies solely on the explicit spectrum of Proposition~\ref{prop:spectrum}.
\end{proof}

\begin{theorem}[Discrete null controllability and convergence of controls]
  \label{thm:discrete-null-control}
  For any $T>0$ and $y_h^0$, there is a control $v_h \in L^2(0,T)$ steering
  the solution of \eqref{eq:discrete-internal} (or its boundary analogue) to
  $y_j(T)=0$, $j=1,\dots,N$. Moreover, for $y^0\in L^2(0,L)$, the controls
  $v_h$ may be built so that
  \begin{equation}
    \label{eq:convergence-controls}
    v_h \to v \ \text{in } L^2(0,T) \ \text{as } h \to 0,
  \end{equation}
  where $v$ is a null control for the continuous heat equation, provided the
  initial data are chosen appropriately.
\end{theorem}
\begin{proof}[Source]
  \citet{lopez1998some}.
\end{proof}

\begin{remark}[Extension to internal control]
  \label{rem:internal-extension}
  Theorem~\ref{thm:discrete-null-control} -- stated by \citeauthor{lopez1998some}
  for the boundary case -- can be modified to also give the convergence of the
  discretization for internal control, system \eqref{eq:discrete-internal}
  \citep[Rem.~3.4]{zuazua2006control}.
\end{remark}

\begin{remark}[Scope: only in 1D]
  \label{rem:1d-only}
  The uniform property of Definition~\ref{def:uniform-controllability} holds
  \emph{only for the one-dimensional} heat equation \citep{lopez1998some};
  convergence \eqref{eq:convergence-controls} follows from an explicit choice
  of $y^0$ via Fourier series \citep[Rem.~1.2]{lopez1998some}; and $v_h$
  converges to the minimal-norm HUM control of
  Section~\ref{ssec:ode-prototype}. In two dimensions, the space semi-discrete
  heat equation is \textbf{not} controllable at all \citep{zuazua2005propagation}
  -- a counterexample that is dimension-general, not specific to internal or
  boundary control, and that directly explains the 1D restriction here.
  Further references establishing/refining uniform controllability for the
  space-discrete 1D case, not worked through in detail here, include
  \citet{labbe2006uniform} and \citet{boyer2010discrete}.
\end{remark}

\begin{remark}[Gap, not filled: joint $(h,k)$ scaling]
  \label{rem:hk-gap}
  Every uniform-controllability result above is for \emph{space}
  semi-discretization only ($h\to0$, time kept continuous); no joint
  $(h,k)$-refinement (fully discrete) uniform-controllability theorem is
  established here. This is the discrete-side counterpart of the open
  condition-number scaling question already noted after
  Proposition~\ref{prop:condition-number}: the penalized-HUM condition number
  $\nu_a = 1+\|\Lambda\|/\varepsilon$ is stated purely as a function of
  $\varepsilon$, with no established $(h,k)$-refinement scaling law. Both gaps
  should be read together, and neither should be presented as resolved.
\end{remark}

\subsection{Other formulations (brief, not pursued here)}
\label{ssec:other-formulations}

The time-discrete and fully-discrete formulations are noted only for
completeness, since this work uses space discretization exclusively.
\textbf{Time-discrete}: this is a genuine \emph{negative} result, not a milder
version of the space-discrete case above -- the time-semi-discrete
multi-dimensional heat equation is in general neither null nor approximately
controllable \citep{zheng2008controllability}, and Carleman estimates for the
time-discrete operator give only "another sense of controllability"
\citep{boyer2020carleman}, not the same uniform-observability property.
\textbf{Fully discrete}: results are scarce; null controllability for a
family of one-dimensional parabolic equations has been shown via Carleman
inequalities, again only in a relaxed sense of controllability
\citep{casanova2021carleman}.

\section{Known theoretical and numerical issues}
\label{sec:known-issues}

\subsection{Internal control}

\begin{remark}[Ill-posedness of the unpenalized minimization]
  \label{rem:coercivity-space}
  The space on which $J$ (Proposition~\ref{prop:pde-gramian}) is coercive is
  much larger than $L^2(0,1)$: it is the completion of $L^2(0,1)$ with respect
  to the norm
  $\big(\tfrac12\int_0^1\!\int_0^T|\mathbf{1}_\omega\varphi|^2\,dx\,dt\big)^{1/2}$
  -- exactly the space $F$ of Remark~\ref{rem:hum-existence}. This space can
  hardly be approximated by a finite-dimensional basis
  \citep{munch2010numerical}, which is why the numerical problem of minimizing
  the dual functional is intrinsically ill-posed.
\end{remark}

Two genuine cost-of-control results appear in the source material: an
exponential-in-\emph{time} estimate,
$\|v\|_{L^\infty(\Sigma)} \le e^{-2\lambda_1 T}C(T-\tau)\|z_0\|_{L^2(\Omega)}$
(here $\lambda_1$ is the first Dirichlet-Laplacian eigenvalue on $\Omega$,
$z_0$ the shifted initial datum, and $\tau$ an intermediate time -- all from
the proof of Theorem~\ref{thm:constrained-linear} in
Section~\ref{sec:constrained-aside}, cost in $T$, not in $\varepsilon$), and a
purely \emph{experimental} control-norm blow-up as the control region $\omega$
shrinks. No exponential-in-$\varepsilon$ blow-up result
(of the form $k\sim e^{1/\varepsilon}$, sometimes seen in the wider numerical
HUM literature) is established in the source material used here, and none is
claimed.

\begin{remark}[Gramian conditioning vs.\ control-region geometry]
  \label{rem:gramian-geometry}
  Although the literature guarantees exact/null controllability for
  internal control on any open $\omega$, with no restriction on its geometry
  \citep{zuazua2002controllability}, the condition number of the Gramian
  $\Lambda$ depends strongly on the location of $\omega$; when discretized in
  space, $\Lambda_N$, it blows up as $N$ increases and $\omega$ shrinks
  \citep{munch2011optimal}.
\end{remark}

\begin{remark}[Time-integrator sensitivity]
  Convergence of the numerical HUM method is highly sensitive to the choice
  of time integrator; explicit schemes blow up for most mesh-parameter pairs,
  even when the underlying finite-difference scheme is convergent and
  uniformly controllable (Section~\ref{sec:discrete}). Gradient-type
  optimization methods can likewise fail even under those same favorable
  conditions; preconditioning or other iterative methods for the resulting
  large linear systems are a natural, and currently active, avenue for
  improving conjugate-gradient convergence (Theorem~\ref{thm:cg-rate}).
\end{remark}

\begin{remark}[Regularity of the adjoint state and the penalization norm]
  There is a strong dependency between the regularity of the adjoint state
  and the minimizer of the dual functional: different choices of penalization
  norm produce genuinely different solution behavior
  \citep{kindermann1999convergence}.
\end{remark}

\subsection{Boundary control}

Boundary control requires a change of norms and function spaces relative to
internal control (Remark~\ref{rem:boundary-h10}), which in practice makes the
numerical method more unstable. For a uniform control-point distribution, a
denser distribution near the endpoints of the control interval may improve
results \citep{kalimeris2021numerical}.

\begin{remark}[Lack of robustness]
  Small perturbations in input data can produce widely different numerical
  results -- a general robustness concern for both control types, not
  specific to either.
\end{remark}


\bibliographystyle{abbrvnat}
\bibliography{theory}

\end{document}
